% called by main.tex
%
\chapter{Preliminary study}
\label{ch::chapter2}
\section{Introduction} 

The first chapter marks the beginning of our project's study, during which we will describe the project, its field, scope, and significance. We will start by introducing the firm that is the center of our discussion. After the review, we conduct a critical analysis of the current situation to find the problems which our Optical Character Recognition (OCR) website aims to solve and the goals it wants to accomplish, along with defining the methodology chosen to deliver this project.

\section{General overview  }
The current project "Optical Character Recognition Application" is done as a part of the End-of-Studies Project for the academic year 2023-2024. It has been done at the company Elite Consulting. 

\section{Company presentation}
In this section, we will provide an overview of the host organization, Elite Council Consulting. We will cover its background, key details, activity areas, and organizational structure.
\subsection{About the host organization}
Elite Council Consulting is a consulting firm located in Hammamet, specializing in helping businesses overcome risks and challenges while guiding them towards the best prospects. The company collaborates with major technology clients worldwide to create new initiatives and increase profits. Established in 2019, it has between 20 and 100 employees.\\
The figure 1.1 below shows the logo of Elite Council Consulting :
\begin{figure}[!htpb]
    \centering
    \includegraphics[height=2cm]{figures/ECC Logo.png} 
    \caption{Elite Council Consulting logo.}
    \label{fig:ecc_logo}
\end{figure}
\subsection{Descriptive sheet}
\begin{itemize}
    \item \textbf{Company name :} Elite Council Consulting
    \item \textbf{Date of establishment :} December 2019
    \item \textbf{Legal Form :} Limited Liability Company
    \item \textbf{Address :} Rue de Nevers 8050 Hammamet
    \item \textbf{Phone number :} 97 380 712
    \item \textbf{Email :} elite.council.consulting@gmail.com
\end{itemize}
\subsection{Activity areas}
The consultants at Elite Council Consulting possess technical skills, business knowledge, and substantial consulting experience. With these assets, the company offers a wide range of consulting services in various fields:

\begin{itemize}
    \item Finance
    \item Marketing
    \item Information Technology
\end{itemize}

\subsection{Organization chart}
An organization chart is a visual representation of a company's structure, showing the relationships and ranks of its parts and positions.
the figure 1.2 below shows the Elite Council Consulting firm organization chart.
\begin{figure}[!htpb]
    \centering
    \includegraphics[height=3cm]{figures/Organization Chart.png} 
    \caption{Organization chart.}
    \label{fig:Elite_Council_Consulting_organization_chart}
\end{figure}

\section{History and definition of Optical Character Recognition}
OCR is a well-known concept that was born in the mid-20th century when the necessity of the process of converting printed or handwritten text into digital formats was felt. OCR technology developed alternatively with the growth of computing and digital imaging. Its main purpose is to extract text data from images, scanned documents, or other media, thus, it allows the machines to understand and process the textual information. Basically, OCR technology changes the way we deal with textual content, making the connection between the physical and digital documents possible, and at the same time, providing  efficiency, accuracy and accessibility advantages 


\section{Existing OCR solutions study}
OCR was designed to convert the documents from different types such as scanned paper and pictures  into a text format, which allows for editing and exporting. The purpose is to improve  document management processes and enhance data accuracy, and improve . Here is an overview of how OCR is currently implemented in the industry: 

1. Adobe Scan

\begin{itemize}
    \item Functionality: Adobe Scan helps users to scan documents and convert their contents into PDFs ;

\end{itemize}

2. Google Drive OCR

\begin{itemize}
    \item Functionality: Google Drive’s OCR feature converts images and PDFs into editable Google Docs;  

\end{itemize}

 3. ABBYY FineReader

\begin{itemize}
    \item Functionality: ABBYY FineReader is a professional OCR software that supports various editable formats and offers robust editing tools.

\end{itemize}
 

\section {Problem statement }
Despite the advancements in the existing OCR technology, there are some limitations that make them ineffective in certain situations. Below are the primary limitations observed in some of the current leading OCR solutions that we have already mentioned in the previous section.

1. Adobe Scan
\begin{itemize}
    \item Limitations: Adobe Scan lacks a phase where users can edit and cut off OCR errors before finally saving the document and allows users to convert into PDFs only;
\end{itemize}        

2. Google Drive OCR
\begin{itemize}
    \item Limitations: Google Drive OCR processes documents directly, often leaving OCR errors uncorrected. It is primarily integrated within the Google ecosystem, limiting export options to other formats;
\end{itemize}

 3. ABBYY FineReader
 \begin{itemize}
     \item Limitations: ABBYY FineReader is a desktop-based application with a steep learning curve, making it less accessible for casual users. It requires significant setup and is not as user-friendly as web-based solutions which can be overwhelming, especially for those who need a simple and quick solution. Additionally, it does not facilitate easy collaboration due to its lack of web-based accessibility. Furthermore, it is a premium software, which might be cost-prohibitive for smaller businesses or individual users. 
\end{itemize}

\section{Proposed solution   }
In order to overcome the limitations of the current OCR systems, our proposed solution provides a more comprehensive and user-friendly method. Our project provides an editing stage during which users check and correct any OCR errors before finalizing the document to increase the accuracy. Also, the users can export the corrected text into formats like Comma-Separated Values(CSV),JavaScript Object Notation (JSON) or Word which bring them more options of the integration. With a user-friendly web interface and our combination of all of these tools our solution makes it easy to use even for those who are not technically advanced.
The figure 1.3 below represents the flowchart Infographic of our OCR Project.
%l text hedha li ken mawjoud fou9 l flowchart infographic, haw ma7tout fi commentaire ken t7eb a9rah ken l9it fih 7aja tnajm testa3melha fel fa9ra hedhi bch twali relevant akthr lel taswira sinon fas5ou :
%This process is divided into six main phases: image upload, OCR scanning, data verification and editing, ETL , data storage and data export. All the phases are designed to work together smoothly,  making sure that data is extracted and managed efficiently.\\
%The first phase involves easily uploading images, which are then processed using OCR technology to identify and extract information. Next, users can review and edit the extracted data to ensure everything is accurate before it moves to the ETL process, where the data is organized into a coherent dataset. After this, the data is stored and made available for flexible export options, as clarified in the project's flowchart infographic figure below.

\begin{figure}[!htpb]
    \centering
    \includegraphics[height=10.13cm]{figures/Flowchart Infographic.jpg}
    \caption{InvoiceScan+ flowchart infographic.}
    \label{fig:flowchart_infographic}
\end{figure}
\newpage


\section{The adopted methodology}
Choosing the right methodology is crucial for the success of any project. It provides an efficient and structured way of achieving objectives and allows the management of deadlines and the quality control of output. 
\subsection{The choice of agile methodology}
agile methodology is a flexible and collaborative approach to project management and development. It focuses on iterative progress, where small parts of the project are completed in cycles, allowing for continuous feedback and improvement. One popular Agile framework is SCRUM, which involves breaking down the project into manageable tasks called sprints. Teams work together closely, adapting to changes and ensuring that the final product meets the evolving needs of the users. This approach emphasizes transparency, customer involvement, and frequent delivery of functional components.
\subsection{Why SCRUM }
SCRUM offers several advantages that make it the optimal choice for our project.
\begin{itemize}

\item Flexibility: SCRUM make room for flexibility in responding to changes in the requirements and priorities. It accommodates iterative development, making adjustments can be made throughout the project life-cycle;

\item Transparency: SCRUM promotes transparency by providing clear visibility into the progress of the project. Through regular meetings and updates, supervisors are kept informed about the project's status and can provide feedback as needed;

\item Team Collaboration: SCRUM helps collaboration among team members by promoting multitask teamwork. Through daily meetings and collaborative decision-making, teams can work together effectively to achieve project goals.
\end{itemize}
The figure 1.4 is a visual representation of how SCRUM works .

\begin{figure}[!htpb]
    \centering
    \includegraphics[height=6cm]{figures/Scrum-figure.jpg}
    \caption{Scrum infographic.}
    \label{fig:Scrum}
\end{figure}


\subsection{Project management with SCRUM}
To understand how SCRUM will be implemented in our project, it is crucial to define its key components and roles. Here's an overview of the basic elements of SCRUM and their application in our project.


\subsubsection{Project planning with SCRUM}
   •Product Backlog: This is an estimated list of all the features, enhancements, and bug fixes required for the project, expressed as "user stories". It serves as the project's to-do list and is continuously updated based on feedback and changing requirements.


   •Sprint Backlog: A subset of the product backlog, the sprint backlog includes the tasks and user stories that the team aims to complete during a specific sprint, typically lasting 2-4 weeks. It helps in planning and tracking the sprint's progress.


   •Increment: At the end of each sprint, the team delivers a potentially shippable product increment. This increment is tested and reviewed to ensure it meets the project's quality standards and can be delivered to the users.


  •SCRUM Meetings: Regular meetings are held to keep the team aligned and ensure smooth progress. 

  
\subsubsection{Team and roles}
A successful SCRUM implementation relies on clearly defined roles within the team. The table  1.1 represents the key roles and their responsibilities.






\begin{table}[!htpb]
    \caption{Roles and responsibilities in SCRUM.}
    \centering
    \begin{tabular}{|p{3cm}|p{6cm}|p{3cm}|}
        \hline
        \textbf{Role} &\textbf{ Mission} & \textbf{Actor} \\ \hline
        Product Owner & Represents stakeholders and the customer, prioritizes the product backlog and ensures the team works on the most valuable features first. & Elite Council Consulting \\ \hline
        SCRUM Master & Facilitates SCRUM process,ensures the team follows SCRUM practices, removes impediments and supports the team in working autonomously. & Mrs Rania mkhinini \\ \hline
        Development Team &  Builds the product increment, self-organizes, has all necessary skills  and works collaboratively without internal hierarchy. & Seif Eddine Mejrissi and Mohamed Amine Zini \\ \hline
    \end{tabular}
    
    \label{tab:team_roles}
\end{table}



\section{Conclusion }

In this chapter, we have provided an overview of the project's scope and objectives, discussed the current state of OCR solutions, identified existing limitations, and proposed our innovative solution. We also outlined the Agile methodology, specifically SCRUM, that will guide our project management process. Moving forward, we will delve into the architecture and planning phase, where we will define the technical requirements and design necessary to bring our OCR application to life.